\documentclass[10pt,a4paper]{amsart}
\usepackage[margin=19mm]{geometry}
\usepackage{amsmath,amssymb,mathtools}
\usepackage[scr=rsfs,cal=euler]{mathalfa}
\usepackage{xfrac}
\usepackage[unicode,hidelinks,bookmarks=false]{hyperref}
\newcommand{\ie}{i.e.\ }
\newcommand{\cf}{cf.\ }
\newcommand{\resp}{respectively\ }
\newcommand{\Subsection}{\subsection}
\providecommand{\nobreakdash}{\nobreak\hbox{-}}
\setlength{\emergencystretch}{2em}
\multlinegap0pt
\usepackage{enumerate}
\usepackage{mathtools}
\usepackage{xcolor}
\usepackage{amssymb}
\usepackage{float}
\newtheorem{theorem}{Theorem}
\title{Sharp Bohr-Type inequalities for certain classes of close-to-convex functions}
\author{Shalini Rana, Naveen Kumar Jain*}
\date{Source: arXiv:2605.22930v1 (CC BY 4.0)}
\begin{document}
\maketitle

\noindent\emph{Source and licence.} Sharp Bohr-Type inequalities for certain classes of close-to-convex functions. Version arXiv:2605.22930v1 (CC BY 4.0), distributed by arXiv under the Creative Commons Attribution 4.0 International licence, http://creativecommons.org/licenses/by/4.0/, as recorded in the arXiv licence metadata for this exact version and shown on the abstract page https://arxiv.org/abs/2605.22930v1. Full licence notice: \texttt{licenses/arxiv-2605.22930-CC-BY-4.0.txt}.\par\smallskip

\noindent\textit{Editorial scope.} Counted unit: the article's theorem th2 (source0.tex lines 805--815). The article's complete original proof of it is delivered with it (source0.tex lines 816--880, one \texttt{proof} environment of 65 lines). No mathematics is written, repaired or re-derived: the statement and the proof are the article's own lines, in the article's own notation, and no retained line is substituted or re-flowed.  Retained uncounted as the source-local context the proof needs: lines 775--777; lines 785--790; lines 792--797.  Nothing was omitted from any retained span, so the package carries no omission record.  No retained line contains a citation, so the wrapper carries no bibliography and no bibliography entry.  No retained line contains a figure environment, an image-inclusion command or drawing code, so no image is delivered and the package declares no asset dependency.  Editorial formatting only, in addition: an AMS \texttt{amsart} preamble that restates the article's own declarations and macros used by the retained lines, namely the environments \texttt{theorem} on separate counters, together with the standard packages those lines need.  The retained environments print in this document as Theorem 1 in order of appearance, whereas the article prints the same items with the numbering of its own full document.  The complete native source of the article as distributed in the arXiv e-print for this exact version is bundled unchanged as \texttt{sources/201169/source0.tex}.
\begin{equation}\label{coef3}
 |a_n| \le \frac{2}{3} + \frac{1}{3n^2}  .
\end{equation}

\begin{equation}
\frac{2r}{3(1+r)}+\frac{1}{3} \int_{0}^{r} \frac{\log(1+t)}{t} dt
\le |f(z)|
\le \frac{2r}{3(1-r)}-\frac{1}{3} \int_{0}^{r} \frac{\log(1-t)}{t} dt
\label{growth3}
\end{equation}

\begin{equation}
\frac{2}{3(1+r)^2}+ \frac{\log(1+r)}{3r}
\le |f^{'}(z)| \le
\frac{2}{3(1-r)^2}- \frac{\log(1-r)}{3r}
\label{distortion3}
\end{equation}

\begin{theorem}\label{th2}
Suppose that $f(z)=z + \sum_{n=2}^{\infty} a_n z^n \in \mathcal{C}_3$, then
\begin{equation}\label{2.3.1}
|f(z)| + |f^{'}(z)||z| + \sum_{n=2}^{\infty} |a_n z^n| \le d( f(0), \partial{f(\mathbb{D})})
\end{equation}
for $|z| \le r_{31}$, where $r_{31}$ is the solution of
\begin{equation*}
\frac{2 (2-r^2) r}{3 (1-r)^2}-\frac{1}{3} \int_{0}^{r} \frac{\log(1-t)}{t} dt-\frac{1}{3} \log(1-r)+ \sum_{n=2}^{\infty} \frac{r^n}{3 n^2}-\left(\frac{1}{3}+\frac{\pi ^2}{36}\right)=0.
\end{equation*}
The radius is sharp.
\end{theorem}

\begin{proof}
We begin the proof by evaluating
\begin{equation*}
d(f(0), \partial{f(\mathbb{D})})= \liminf \limits_{|z| \rightarrow 1} |f(z)-f(0)|.
\end{equation*}
From equation ($\ref{growth3}$), we get
\begin{equation*}
d(f(0), \partial{f(\mathbb{D})}) \ge \lim \limits_{r \rightarrow 1} \left(\frac{2r}{3(1+r)}+\frac{1}{3} \int_{0}^{r} \frac{\log(1+t)}{t} dt \right)
\end{equation*}
which gives
\begin{equation}\label{eq16}
d(f(0), \partial{f(\mathbb{D})}) \ge \frac{1}{3} +\frac{1}{3} \int_{0}^{1} \frac{\log(1+t)}{t} dt=\frac{1}{3}+\frac{\pi ^2}{36}.
\end{equation}
From  equations $(\ref{coef3})$, $(\ref{growth3})$ and $(\ref{distortion3})$, for all $|z| \le r$, we have
\begin{align}\label{eq17}
|f(z)|& + |f^{'}(z)||z| + \sum_{n=2}^{\infty} |a_n| |z|^n\\
&\le \frac{2r}{3(1-r)}-\frac{1}{3} \int_{0}^{r} \frac{\log(1-t)}{t} dt+ \frac{2r}{3(1-r)^2}-\frac{1}{3} \log(1-r)+ \sum_{n=2}^{\infty} \left( \frac{2}{3} + \frac{1}{3 n^2} \right)  r^n\nonumber\\
&=\frac{2 (2-r^2) r}{3 (1-r)^2}-\frac{1}{3} \int_{0}^{r} \frac{\log(1-t)}{t} dt-\frac{1}{3} \log(1-r)+ \sum_{n=2}^{\infty} \frac{r^n}{3 n^2}.
\end{align}
In order to find the radius $r_{31}\in (0, 1)$ satisfying the inequality $(\ref{2.3.1})$, we define the function $P:[0, 1)\rightarrow\mathbb{R}$ by

\begin{equation*}
P(r)=\frac{2 (2-r^2) r}{3 (1-r)^2}-\frac{1}{3} \int_{0}^{r} \frac{\log(1-t)}{t} dt-\frac{1}{3} \log(1-r)+ \sum_{n=2}^{\infty} \frac{r^n}{3 n^2}-\left(\frac{1}{3}+\frac{\pi ^2}{36}\right)
\end{equation*}
An easy calculation shows that
\begin{equation}\label{eq14}
P(0)=-\left(\frac{1}{3}+\frac{\pi ^2}{36}\right)<0
\end{equation}
and
\begin{align}\label{eq15}
P\left(\frac{1}{2}\right)&=2-\frac{\pi ^2}{36}+\frac{\log2}{3}+\frac{1}{36} \left(-6+\pi ^2-6 \log 2^2\right)+\frac{1}{36} \left(\pi ^2-6 \log 2^2\right)\nonumber\\
&=\frac{1}{36} \left(66+\pi ^2-12 \log 2^2+\log4096\right)\nonumber\\
&\simeq2.17839>0.
\end{align}
Equations \eqref{eq14} and \eqref{eq15} together with Intermediate Value Theorem yield that there exists a root $r\in(0, 1/2)$ of the equation $P(r)=0.$
Also,
\begin{align*}
P'(r)&=\frac{2 \left(2+2 r-3 r^2+r^3\right)}{3 (1-r)^3}-\frac{\log(1-r)}{3r}+\frac{1}{3 (1-r)}-\frac{1}{3} \left(1+\frac{\log(1-r)}{r}\right)\\
&=\frac{1}{3} \left(\frac{4-r (-5+(8-3 r) r)}{(1-r)^3}-\frac{2 \log(1-r)}{r}\right)\\
&=\frac{r \left(4+5 r-8 r^2+3 r^3\right)-2 (1-r)^3 \log(1-r)}{3 (1-r)^3 r}\\
&=\frac{r \left(4+(1-r) r (5-3 r)\right)-2 (1-r)^3 \log(1-r)}{3 (1-r)^3 r}\\
&>0.
\end{align*}
Thus, there exists the  unique  root $r\in(0, 1)$ of the equation $P(r)=0$, say $r_{31}$. By equations \eqref{eq16} and \eqref{eq17}, for $|z|\leq r_{31}$, we have
\[|f(z)| + |f^{'}(z)||z| + \sum_{n=2}^{\infty} |a_n z^n| \le d( f(0), \partial{f(\mathbb{D})}).\]
To prove the sharpness, consider the function
\begin{equation*}
    f(z)= \frac{2z}{3(1-z)} - \frac{1}{3} \int_{0}^{z} \frac{\log(1-\zeta)}{\zeta} d\zeta.
\end{equation*}
Note that
\[-\frac{1}{3}\int_0^r \frac{\log(1-t)}{t} \, dt=\frac{r}{3}+\frac{r^2}{12}+\frac{r^3}{27}+\frac{r^4}{48}+\frac{r^5}{75}+\frac{r^6}{108}+\cdots\]
and \[-(1/3) \log(1 - r)=\frac{r}{3}+\frac{r^2}{6}+\frac{r^3}{9}+\frac{r^4}{12}+\frac{r^5}{15}+\frac{r^6}{18}+\cdots\]
So, at  $z=r=r_{31}$, we have
\begin{align*}
|f(z)|& + |f^{'}(z)||z| + \sum_{n=2}^{\infty} |a_n| |z|^n\\
&= \left|\frac{2r}{3(1-r)} - \frac{1}{3} \int_{0}^{r} \frac{\log(1-\zeta)}{\zeta} d\zeta \right| + |r| \left| \frac{2}{3(1-r)^2}-\frac{\log(1-r)}{3r} \right|\\
  &+   \sum_{n=2}^{\infty} \left| \left( \frac{2}{3}+ \frac{n^2}{3}\right) \right| |r|^n\\
&= \frac{2r}{3(1-r)} - \frac{1}{3} \int_{0}^{r} \frac{\log(1-\zeta)}{\zeta} d\zeta  +   \frac{2r}{3(1-r)^2}-\frac{\log(1-r)}{3r}  +   \sum_{n=2}^{\infty}  \left( \frac{2}{3}+ \frac{n^2}{3}\right) r^n\\
&=\frac{2 (2-r^2) r}{3 (1-r)^2}-\frac{1}{3} \int_{0}^{r} \frac{\log(1-t)}{t} dt-\frac{1}{3} \log(1-r)+ \sum_{n=2}^{\infty} \frac{r^n}{3 n^2}\\
&=\frac{1}{3}+\frac{\pi ^2}{36}\\
&=\frac{1}{3} +\frac{1}{3} \int_{0}^{1} \frac{\log(1+t)}{t} dt\\
&=d(f(0), \partial{f(\mathbb{D})}).
\end{align*}
Hence, the result is sharp.
\end{proof}

\end{document}
